\section{Neural Models used for Compression}

We evaluate our compression strategy on two types of neural models that are commonly used for classification tasks:

%To test our compression strategy we test 2 different architectures, the first is the DAN architecture~\cite{iyyer2015deep} and the second is a LSTM based classfication architecture~\cite{joulin2016bag}. 

\subsubsection{LSTM}
Long-Short term memory neural networks (LSTMs)~\cite{hochreiter_lstm}, are powerful and popular models for classification tasks. The LSTM unit is an RNN extension, that is designed to handle long term dependencies through the use of an input gate, an output gate, a forget gate and a memory cell. Given a sentence input $\overline{X}=x_1, \cdots x_T $ and corresponding word embedding input vectors $e_i, i=1, \cdots T$, the LSTM computes a representation $r_t$ at each $t$, which is denoted as $r_t = \phi ( e_t ,r_{t-1} )$. Here, for sentence classification, we collect the full sentence representation obtained from the final timestep of the LSTM and pass it through a softmax for classification:

\begin{gather}
r^{sent} = r^{forw}_{T}, \space  
\hat{S}=softmax( W_{s} r^{sent} + b_{s})
\end{gather}
where $r^{sent}$ is the sentence representation, and $\hat{S}$ is the label predicted for the sentence. For our classification experiments, we use 300-dimensional Glove ~\cite{pennington2014glove} embeddings as input. Our vocabulary size $V$ contains only the words that appear in the training data, and we compress the input word embedding matrix $W$ (original size $V \times 300$).

\subsubsection{DAN}
The Deep Averaging Network (DAN) model proposed ~\cite{iyyer2015deep} is a bag-of-words neural model that averages the word embeddings in each input utterance as the sentence representation $r^{sent}$. $r^{sent}$ is passed through a series of fully connected layers, and fed into a softmax layer for classification. Assume an input sentence  $\overline{X}$ of length $T$, and corresponding word embeddings $e_i$, then the sentence representation is:  
\begin{gather}
r^{sent} = \frac{1}{T} \sum_{i=1}^T { e_i }
\end{gather}
DAN has been proven a state-of-the-art model for simple classification tasks such as the Factoid Question Answering task~\cite{iyyer2015deep}. In out setup, we have two fully connected layers after the sentence representation $r^{sent}$, of sizes 1024 and 512 respectively. For the DAN experiments we tried DAN-RAND \cite{iyyer2015deep} variant where we take a randomly initialized 300 dimensional embedding. For this experiment our Baseline2 is to use appropriate low rank randomly initialized embedding matrix and train on that from the beginning. Setting up the experiments with randomly initialized as well as GloVe initialized ensures that all the gains we report is not due to low dimensional latent structure of the initialization embedding matrix but are attributed to our proposed approach. 
%Similarly to the LSTM experiments, we use 300-dimensional Glove embeddings as input and compress the input word embedding matrix of original size $V \times 300$.

%% To test our compression algorithm we try 2 different settings
%% %% \begin{itemize}                 


%% Our network consists of one hot vector of size |V| as input, 2 fully connected layers as shown in Figure\ref{} for classification. The layers have the  size 1024 and 512 respectively. For this setting  we compress the fully connected layers by replacing them with smaller layers and fine tune them.  We also train a neural network from scratch with the compressed architecture to showcase the gains  come from the SVD compression and not from the  bottleneck introduced in the architecture architecture. The results for this experiment are shown in  Table 1. Since the initial input layer of size |V|x|H|  was the largest layer by a margin, we only apply compression on it. 





%% We perform all our experiments on 1 Tesla K80 gpu, use Adam as the optimizer with learning rate of 0.001. We use dropout of 0.4 between layers and L2 regularization with weight of 0.005. 
%% We plot the validation fine tuned accuracy and compression ratio's in
%% Figure~\ref{plot:sst_lstm} and Figure~\ref{plot:sst_dan}. Results on the test set are shown in Table~\ref{results}. We show the baseline accuracy on top for each experiment (compression ratio 0). We show the initial test accuracy after SVD and final test accuracy after fine tuning.
